\section{GDTR 结构体}

\texttt{lgdt} 指令的操作数是一个 6 字节（48 位）的结构体：

\codefile{src/include/arch/gdt.h}
\begin{lstlisting}[style=cstyle]
typedef struct __attribute__((packed)) {
    uint16_t limit;  // GDT 总字节数 - 1
    uint32_t base;   // GDT 数组的线性地址
} gdtr_t;
\end{lstlisting}

\begin{figure}[H]
  \centering
  % figures/ch03/fig07-gdtr-6.tex — auto-extracted from chapter source
\begin{tikzpicture}[node distance=0cm]
  \node[memcell, minimum width=2.5cm, fill=green!10] (lim) {limit (16 bit)};
  \node[memcell, minimum width=5cm, fill=blue!10, right=0cm of lim] (base) {base (32 bit)};
  \node[memaddr, left=0.3cm of lim] {+0};
  \node[memaddr, left=0.3cm of base] {};

  % Arrow from base to GDT
  \node[below=1.2cm of base, font=\small] (gdt) {gdt[0] | gdt[1] | gdt[2]};
  \draw[pointer] (base.south) -- (gdt.north) node[midway, right, font=\scriptsize] {指向 GDT 数组};
\end{tikzpicture}

  \caption{GDTR 的 6 字节布局}
\end{figure}

\codefile{src/kernel/arch/gdt.c}
\begin{lstlisting}[style=cstyle]
gdtr.limit = (uint16_t)(sizeof(gdt) - 1);  // 3 * 8 - 1 = 23
gdtr.base  = (uint32_t)(uintptr_t)&gdt[0]; // GDT 数组的线性地址
\end{lstlisting}

\begin{thinkbox}
为什么 limit 是 ``sizeof(gdt) - 1`` 而非 ``sizeof(gdt)``？这是 x86 的惯例：
limit 表示最后一个有效字节的\textbf{偏移}，所以 3 个描述符 $\times$ 8 字节 = 24 字节，
最后一个字节偏移是 23。如果写 24 会怎样？
CPU 会认为 GDT 多了一个字节，可能读到 GDT 数组之外的垃圾数据。
\end{thinkbox}

% ============================================================================

